\documentclass[
	12pt,
	openright,
	twoside,
	a4paper,
	oneside,
	brazil,
	%chapter=TITLE,
	%section=TITLE,
	%subsection=TITLE,
	sumario = tradicional
]{abntex2}

\usepackage[font=small,format=plain,labelfont=bf,up,margin=1cm]{caption}
\usepackage{amsmath}
\usepackage[utf8]{inputenc}
\usepackage[T1]{fontenc}
\usepackage[pdftex]{graphicx}
\usepackage{subfig}
\usepackage{float}
\usepackage[alf,abnt-etal-list=0, abnt-etal-cite=2, abnt-etal-text=it, abnt-repeated-author-omit=yes, abnt-emphasize=bf]{abntex2cite}
% alf - exibi as referencias em ordem alfabetica.
% abnt-etal-list=0 - foi utilizado para não permitir o uso de et al. nas referências.
% abnt-etal-cite=2 - mostra que pode ser feita a citação de até 2 autores no texto. Se numa referência existem 3 autores, usa-se evidentemente o et al., com a citação do primeiro autor.
% abnt-etal-text=it - exibi a escrita do et al nas referencias em italico.

\usepackage[table,xcdraw]{xcolor} % pacote para colorir células da tabela
\renewcommand{\arraystretch}{1.2} % espaçamento da linha na tabela

\usepackage{indentfirst} % justificar automaticamente os parágrafos
\usepackage{lscape}
\usepackage{import}
\usepackage{helvet}
\renewcommand{\familydefault}{\sfdefault}
\usepackage{lastpage}			% Usado pela Ficha catalográfica
\usepackage{color}				% Controle das cores
\usepackage[table]{xcolor}      % cores nas células das tabelas
\usepackage{microtype} 			% para melhorias de justificação

\usepackage{lipsum}             % gerador de lero lero

% Acrescentado para gerar as tabelas 22/07/2017
\usepackage[normalem]{ulem}
\useunder{\uline}{\ul}{}

\usepackage{longtable}
\usepackage{empheq}
\usepackage{multirow}
\usepackage{booktabs}
\usepackage{pdfpages} %adicionar arquivos em pdf

\newcommand*\widefbox[1]{\fbox{\hspace{2em}#1\hspace{2em}}}
\newcommand{\head}[1]{\textnormal{\textbf{#1}}}
\newcolumntype{C}[2]{>{\centering\vspace{#2}\let\newline\\\arraybackslash\hspace{0pt}\vspace{#2}}m{#1}}

% Préambulo --------------------------------------------------------------------------------------------------------------------------

\newcommand{\universidade}{Universidade Federal do Vale do São Francisco }

%%%%%%%%%%%%%%%% DADOS PARA CAPA E CONTRA CAPA %%%%%%%%%%%%%%%% 
\newcommand{\curso}{Inserir curso}
\newcommand{\graduacao}{Curso de Graduação em \curso}
\newcommand{\bancaum}{Avaliador 1\\ Faculdade}
\newcommand{\bancadois}{Avaliador 2\\ Faculdade}
\newcommand{\bibliotecario}{Inserir Bibliotecário}
\newcommand{\cpf}{000.000.000-00}
\newcommand{\siglaColegiado}{CENEL}
\newcommand{\siglaUniversidade}{UNIVASF}
\newcommand{\periodo}{2018.1}
\newcommand{\tipodetrabalho}{PROJETO PRELIMINAR}

\autor{Inserir seu nome}

\titulo{Inserir título}

\tipotrabalho{Trabalho de Conclusão de Curso}
\orientador[Orientador: Dr.]{Inserir nome do orientador}
%\coorientador[Coorientador: Dr.]{Inserir nome do coorientador}
\instituicao{\universidade \par \curso}
\local{Juazeiro - BA}
\data{2018}

\preambulo{
	Pré-projeto de Trabalho de Conclusão de Curso apresentado como requisito parcial para obtenção do título de Bacharel em \curso, pela \universidade - UNIVASF.
}

 %%%%%%%%%%%%%%%%%%%%%%%%%%% FIM %%%%%%%%%%%%%%%%%%%%%%%%%%% 

% alterando o aspecto da cor azul.
\definecolor{blue}{RGB}{41,5,195}

% Configurações de aparência do PDF final.
\makeatletter
\hypersetup {
	pdftitle		= {\@title},
	pdfauthor		= {\@author},
	pdfsubject      = {\imprimirpreambulo},
	pdfcreator		= {LaTeX with abnTeX2},
	pdfkeywords 	= {Geração Distribuída}{Biogás}{Biometano}{Energia Renovável}{Aterro Sanitário}{Créditos de Carbono}{Geração de Energia},
	colorlinks		= true,			% false: boxed links; true: colored links
	linkcolor		= blue,         % cor links internos.
	citecolor		= blue,       	% cor links referencias bibliograficas.
	filecolor		= magenta,		% cor links de arquivos.
	urlcolor		= blue,			% Cor url.
	bookmarksdepth	= 4
}

\makeatother

% compila o sumário
\makeindex

% Altera as margens padrões
\pagenumbering{arabic}

\setlrmarginsandblock{3cm}{2cm}{*}
\setulmarginsandblock{3cm}{2cm}{*}
\checkandfixthelayout

% Espaçamentos  entre linhas e parágrafos 
% O tamanho do parágrafo é dado por:
\setlength{\parindent}{1.5cm}
% Controle do espaçamento entre um parágrafo e outro:
\setlength{\parskip}{0.2cm}  % tente também \onelineskip

% Retira espaço extra obsoleto entre as frases.
\frenchspacing

\begin{document}

	\selectlanguage{brazil}

	% ELEMENTOS PRE-TEXTUAIS

	\pretextual

	% Capa
	
	\begin{capa}
		\center

		\vspace*{-2.5cm}

		\begin{figure}[!htbp]
		 	\centering
		 	\includegraphics[scale=0.08]{./figuras/univasf.eps}
		\end{figure}

		\ABNTEXchapterfont {\large \textbf {\MakeUppercase{\universidade}\\ \MakeUppercase{\graduacao} }}
		\vspace{1cm}
		\begin{center}
			\ABNTEXchapterfont \bfseries {\Large \tipodetrabalho}
		\end{center}

		\vspace*{1.5cm}

		\ABNTEXchapterfont {\Large \imprimirautor}

		\vfill
		\begin{center}
			\ABNTEXchapterfont \bfseries {\Large \imprimirtitulo}
		\end{center}	
		

		\vfill

		{\large \imprimirlocal} \\
		{\large \imprimirdata}

		\vspace*{1cm}
	\end{capa}

	
	% Folha de Rosto
	\makeatletter

\begin{center}
		\ABNTEXchapterfont {\large \textbf {\MakeUppercase{\universidade}\\ \MakeUppercase{\graduacao} }}
		\vspace*
		\fill

		{\ABNTEXchapterfont\Large\imprimirautor}

		\vspace*{\fill}

		\begin{center}
			{\ABNTEXchapterfont\bfseries\Large\imprimirtitulo}
		\end{center}

		\vspace*{\fill}\vspace*{\fill}

		\abntex@ifnotempty{\imprimirpreambulo} {
			\hspace{.45\textwidth}
			\begin{minipage}{.5\textwidth}
				\SingleSpacing
				\imprimirpreambulo \\\\
				\imprimirorientadorRotulo~\imprimirorientador\\
				\imprimircoorientadorRotulo~\imprimircoorientador
			\end{minipage}
			\vspace*{\fill}
		}

		\vspace*{\fill}\vspace*{\fill}

		{\large\imprimirlocal}
		\par
		{\large\imprimirdata}
\end{center}

\newpage
\makeatother
	
	% Termo de conformidade
	\makeatletter
\begin{center}
		\ABNTEXchapterfont {\large \textbf {\MakeUppercase{\universidade}\\ \MakeUppercase{\graduacao} }}
		\vspace*
		\fill

		\begin{center}
			{\ABNTEXchapterfont\bfseries\Large\imprimirtitulo}
		\end{center}

		\vspace*{\fill}

		{\Large{PROJETO PRELIMINAR}}

		\vspace*{\fill}\vspace*{\fill}
		
		\assinatura{\imprimirautor\\ Orientado(a)}

		\assinatura{\imprimirorientador\\ Orientador(a)}

		\vspace*{\fill}\vspace*{\fill}

		{\large\imprimirlocal}
		\par
		{\large\imprimirdata}
	\end{center}

	\newpage
\makeatother
	
	% Declaração de orientação
	\newpage

\vspace*{-2.5cm}

\begin{figure}[!htbp]
	\centering
	\includegraphics[width=3.25cm]{./figuras/brasao.eps}\\
	\small{\textbf{SERVIÇO PÚBLICO FEDERAL}}\\
	\textbf{MINISTÉRIO DA EDUCAÇÃO}\\
	\textbf{SECRETARIA DE EDUCAÇÃO SUPERIOR}\\
	\MakeUppercase{\textbf{\universidade}}
\end{figure}

\vspace{2cm}

\begin{center}
	\textbf{\Large{\MakeUppercase{Declaração de Orientação}}}
\end{center}

\vspace{2cm}

\noindent{Declaro para os devidos fins minha aceitação como orientador do trabalho de conclusão de curso (TCC), do discente \textbf{\imprimirautor}, intitulado “\imprimirtitulo”, durante o período letivo \periodo.}
	
\noindent Por ser verdade, firmo o presente.

\vspace{1cm}

\noindent\imprimirlocal, \today.


\begin{center}
	\assinatura{\textbf{\emph{Prof. \imprimirorientador}}\\\textbf{\siglaColegiado/\siglaUniversidade}}
\end{center}

\newpage
	
	% Lista de Figuras
	
	\renewcommand*\listfigurename{Lista de Figuras}
	\listoffigures*
	\cleardoublepage	
	
	% Lista de Tabelas
	\renewcommand*\listtablename{Lista de Tabelas}
	\listoftables*
	\cleardoublepage

	% Índice (Sumário)
	\tableofcontents*
	\cleardoublepage
	
	%%%%%%%%%%%%%%%%%%% ELEMENTOS TEXTUAIS %%%%%%%%%%%%%%%%
	
	%\begin{resumo}

Escrever o resumo aqui\\\\

\vspace{\onelineskip}
\noindent\textbf{Palavras-chaves}: Gerenciamento. Energia Elétrica. Comunicação \textit{WiFi}. \textit{Mobile}.
\end{resumo}

	% --- resumo em inglês ---
\begin{resumo}[Abstract]
\begin{otherlanguage*}{english}

Write here\\\\
\vspace{\onelineskip}
\noindent\textbf{Keys}: Management. Electricity. Wireless communication. Mobile.

\end{otherlanguage*}
\end{resumo} %retire o comentário para inserir
	
	\setcounter{page}{1}
	
	\textual
	% Arquivos para compilação separada
	\chapter{Introdução}

Aqui está um exemplo de como utilizar uma referência  bibliográfica que existe no arquivo \textit{bibitex.bib}, localizado na pasta "bibliografia" \cite{atlas:2006}. Segundo \citeonline{toledo:2010}, testando aqui outra referencia.

\lipsum[4-5]


    \chapter{Justificativa}


Segundo \citeonline[pág. 12]{atlas:2006} \lipsum[1-2]
    \chapter{Objetivos}

\lipsum[1]

\section{Objetivos Específicos}

\begin{itemize}
    \item Lorem ipsum dolor sit amet, consectetur adipisicing elit, sed do eiusmod tempor incididunt ut labore et dolore magna aliqua;
    \item Lorem ipsum dolor sit amet, consectetur adipisicing elit, sed do eiusmod tempor incididunt ut labore et dolore magna aliqua;
    \item Lorem ipsum dolor sit amet, consectetur adipisicing elit, sed do eiusmod tempor incididunt ut labore et dolore magna aliqua.
\end{itemize}

\section{Organização do Trabalho}

A construção deste trabalho de conclusão de curso foi estruturada e subdividida em mais quatro capítulos, além deste introdutório. São eles:

\begin{itemize}
    \item O Capítulo 1 apresenta inicialmente uma visão geral da Lorem ipsum dolor sit amet, consectetur adipisicing elit, sed do eiusmod tempor incididunt ut labore et dolore magna aliqua \textit{on-grid}.
    
    \item O Capítulo 2 trata do uso de \textit{softwares } pLorem ipsum dolor sit amet, consectetur adipisicing elit, sed do eiusmod tempor incididunt ut labore et dolore magna aliqua
    
    \item O Capítulo 3 apresenta o DLorem ipsum dolor sit amet, consectetur adipisicing elit, sed do eiusmod tempor incididunt ut labore et dolore magna aliqua
    
    \item Por fim, o Capítulo 5 expõe as conclusões finais obtidas a partir do estudo e realização deste trabalho.
\end{itemize}
    \chapter{Metologia da Pesquisa}

\lipsum[2-3]

	
	%%%%%%%%%%% ELEMENTOS POS-TEXTUAIS %%%%%%%%%%%
	\postextual
	% Considerações finais
	% \thispagestyle{empty}

	% Referências
	\chapterstyle{abnt}
	\bibliographystyle{abntex2-alf}
	\bibliography{bibliografia}

	%Apêndice
	%\appendix
	 
	%%%%%%%%%%%%%%%%%% FIM %%%%%%%%%%%%%%%%
	 
	%%%%%%%%%%% INDICE REMISSIVO %%%%%%%%%%%
    \phantompart
    \printindex
    %%%%%%%%%%%%%%%%%%% FIM %%%%%%%%%%%%%%%%

\end{document}
